\clearpage
\begingroup
\fancyhead[L]{\small BEMERKUNGEN}
\renewcommand{\headrulewidth}{0pt}
\par\noindent\hspace*{\parindent}Zu [VII.]: Die Briefstellen sind nach den Originalen abgedruckt.
Notiz [2.] stammt vollständig vom Bearbeiter. Die außerordentlich umfangreichen
Rechnungen zu diesen Tafeln (vgl. S. 429) und die fertige Tafel selbst bilden
ein Konvolut aus sorgfältig geordneten Blättern. Viele davon stammen
aus anderer Hand als der von \textsc{Gauss}, nämlich von \textsc{Encke} und \textsc{Westphal}.

\par\noindent\hspace*{\parindent}Zu [VIII.]: Über die Störungen der Pallas durch Saturn enthält der
Nachlass nichts von \textsc{Gauss}' Hand. Er enthält nur die Originalbriefe von
\textsc{Nicolai} an \textsc{Gauss} sowie einige Blätter von \textsc{Nicolai}s Hand.
Diesen Blättern entstammen die Gleichungen 1) und 2) der Notiz [1.]. Der
übrige Inhalt dieser Notiz wurde ergänzt. Die betreffenden Briefstellen
sind in den Notizen [2.]--[6.] abgedruckt. Die Ergebnisse der Störungsrechnung stehen auf
einzelnen Blättern, die ursprünglich den Briefen \textsc{Nicolai}s beilagen
und jetzt im Nachlass in die Kapseln mit den Pallasstörungen eingeordnet sind.
Diese Ergebnisse sind in etwas anderer Reihenfolge abgedruckt als
auf den Originalblättern. Dort sind sie in Manipel geordnet. Für den
Abdruck wurde ihre ursprüngliche natürliche Reihenfolge gewählt.

\par\noindent\hspace*{\parindent}Die ganze Rechnung dieser Störungen ist in \textsc{Nicolai}s Nachlass
an der Großherzoglichen Sternwarte in Heidelberg erhalten. Herr
\textsc{Valentiner} hat sie dem Unterzeichneten freundlicherweise zur
Einsicht überlassen. Sie trägt die Unterschrift »Geendigt den
20\textsuperscript{ten} April 1814.« Vielleicht dürfen wir auf ihre etwas ausführlichere
Veröffentlichung von anderer Seite hoffen. In \textsc{Gauss}' Werken wäre sie wohl kaum angebracht gewesen.
Übrigens hat \textsc{Nicolai} zwei Rechnungen der
Saturnstörungen ausgeführt. Die erste bezog sich nur auf die Hauptgleichungen.
Ihre Ergebnisse sind ebenfalls im Nachlass vorhanden. Die entsprechenden
Briefe von \textsc{Gauss} an \textsc{Nicolai} scheinen verloren zu sein. Jedenfalls
teilte mir Herr \textsc{Valentiner} mit, dass sich in \textsc{Nicolai}s Nachlass nur
eine Reihe von Briefen von \textsc{Gauss} aus späterer Zeit (von 1819 an) findet,
die Herr \textsc{Valentiner} herausgegeben hat (Briefe von C. F.
\textsc{Gauss} an B. \textsc{Nicolai}, Karlsruhe 1877). Die Briefe von \textsc{Nicolai}
an \textsc{Gauss} scheinen vollständig erhalten zu sein. Ihre Veröffentlichung
steht noch aus.

\par\noindent\hspace*{\parindent}Zu [IX.]: Die Formeln und Zahlen der Notiz [1.] finden sich teils im
Handbuch Bc, teils auf Blättern, teils in einem Heft mit dem Titel »Astronomische Rechnungen;
Angefangen im März 1814«, das sonst andere Gegenstände enthält
und daher im Nachlass unter »Beobachtungen und Rechnungen« eingeordnet ist.
Die Zahlen der Notizen [2.]--[3.] stammen aus Rechnungen, die \textsc{Gauss} auf einzelnen Blättern
führte und nicht abschloss. Die Notiz [4.]
stammt vom Bearbeiter. [5.]--[6.] stammen aus dem erwähnten Heft »Astronomische Rechnungen«,
einzelne Beziehungen der Notiz [5.] auch von einem losen Blatt. [7.] stammt schließlich
aus dem Handbuch Be, wo auf den abgedruckten Text noch ein Rechenbeispiel
(für \(M=165^0\)) folgt.

\par\noindent\hspace*{\parindent}Die Notizen [4.]--[7.] betreffen einen sehr interessanten
Kunstgriff zur Entwicklung der Störungsfunktion. \textsc{Gauss} hat ihn
schließlich aufgegeben, weil er offenbar keine wesentliche Einsparung
beim numerischen Rechnen bringt. Nach dieser
Methode hat \textsc{Gauss} außer den abgedruckten Werten für \(M=30^0\) nur noch
die für \(M=15^0\) und \(45^0\) berechnet.

\medskip
\par\noindent\hspace*{\parindent}Man sieht, dass \textsc{Gauss} das gewaltige Problem, die
Pallasstörungen innerhalb der Genauigkeitsgrenze der Beobachtungen zu berechnen, fast ganz
gelöst hat. Auch heute würde sich ein Astronom nur ungern daran wagen.
Es fehlen nur wenige Störungsgleichungen für den Mars.
Umso mehr war zu bedauern, dass diese Arbeit bis heute
unbekannt geblieben ist. Wie aus den letzten abgedruckten Briefstellen
hervorgeht, hat auch \textsc{Gauss} selbst dies lebhaft empfunden.

\par\noindent\hspace*{\parindent}Sehr auffällig ist, dass \textsc{Gauss} in seinem wissenschaftlichen
Tagebuch (oder Notizenjournal, wie er es selbst nennt), das von Herrn
F. \textsc{Klein} in der »Festschrift zur Feier des 150jährigen Bestehens der K. Ge%
\SourceLeaf{616}sellschaft der Wissenschaften zu Göttingen«, 1901 herausgegeben
worden ist, diese Untersuchung kein einziges Mal erwähnt.
Über die Jahre 1802--1804 steht im Tagebuch die Bemerkung: »\textit{Annis
insequentibus 1802, 1803, 1804 occupationes astronomicae maximam otii partem
abstulerunt, calculi imprimis circa planetarum novorum theoriam instituti.
Unde evenit, quod hisce annis catalogus hicce neglectus est} \mbox{.\ .\ .\ .\ .\ .\ .}«
In diese Jahre fallen \textsc{Gauss}' Arbeiten zum Zweikörperproblem
(zur Vorbereitung der Theoria motus) und zu den Störungen der Ceres. Merkwürdig
ist, dass \textsc{Gauss}' Beschäftigung mit astronomischen Rechnungen
stets dazu führte, dass er die Führung des Tagebuchs unterbrach. Denn auch
im Jahre 1810, als \textsc{Gauss} begann, sich mit der Pallas zu beschäftigen,
beginnt eine große Lücke. Am 29. Februar 1812 schreibt er dann: »\textit{Catalogum
praecedentem per fata iniqua iterum interruptum initio anni 1812 resumimus}
\mbox{.\ .\ .\ .\ .}« Er erwähnt dann, dass er den »zweiten« Beweis
des Fundamentalsatzes der Algebra wiedergefunden habe. Er hatte ihn im November 1811 gefunden,
dann aber zum Teil wieder vergessen, da er nichts
aufgezeichnet hatte. Zwischen dieser und der nächsten Notiz vom 26. September
1812, die sich auf die Anziehung eines Ellipsoids bezieht, liegt die
Entdeckung der Rationalität der mittleren Bewegungen von Pallas und Jupiter.
Auch über sie schweigt das Tagebuch vollständig.

\par\noindent\hspace*{\parindent}Dass die \textsc{Gauss}schen Untersuchungen über die Störungen der Pallas
noch heute wertvolles Material liefern, liegt auf der Hand. Die Reduktion
der Störungsausdrücke auf die heutige Epoche ist
sehr wünschenswert.
\par\nobreak\noindent\hfill\textsc{Brendel}.\par
\nobreak\vspace{3pt}\noindent\makebox[\linewidth]{\rule{.24\linewidth}{.4pt}}\par
\clearpage
\endgroup
